%=============================================================================
% 模块四：CV 领域深度实战——阅卷系统的第一步 (约 12 页)
%=============================================================================

\section{CV领域深度实战}

% -----------------------------------------------------------------------------
% 1. 任务背景
% -----------------------------------------------------------------------------

\begin{frame}{任务背景：试卷图像的噪声处理}
	\begin{block}{故事问题}
		手机拍摄的试卷照片总是有噪点、光照不均,如何自动处理这些问题？
	\end{block}

	\vspace{0.3cm}

	\begin{columns}
		\column{0.5\textwidth}
		\textbf{常见图像质量问题：}
		\begin{itemize}
			\item \textbf{噪声：} 手机传感器产生的随机噪点
			\item \textbf{光照不均：} 拍摄环境光线不均匀
			\item \textbf{模糊：} 手抖或对焦不准
			\item \textbf{透视变形：} 拍摄角度倾斜
			\item \textbf{阴影：} 遮挡造成的暗区
		\end{itemize}

		\column{0.5\textwidth}
		\textbf{对阅卷系统的影响：}
		\begin{itemize}
			\item 二值化效果差 $\rightarrow$ 填涂识别错误
			\item 边缘检测不准 $\rightarrow$ 答题卡定位失败
			\item 文字识别率低 $\rightarrow$ 主观题评分困难
		\end{itemize}

		\vspace{0.3cm}

		\begin{exampleblock}{本周目标}
			用 \textbf{AI 辅助} 实现：
			\begin{enumerate}
				\item 图像去噪（高斯/中值滤波）
				\item 透视矫正（透视变换）
				\item 边缘检测优化
			\end{enumerate}
		\end{exampleblock}
	\end{columns}
\end{frame}

% -----------------------------------------------------------------------------
% 2. 实战 1：引导 AI 写出滤波代码
% -----------------------------------------------------------------------------

\begin{frame}{实战 1：引导 AI 写出滤波代码}
	\tiny
	\vspace{-0.8cm}
	\textbf{任务：}对比高斯滤波、中值滤波效果

	\begin{columns}
		\column{0.5\textwidth}
		\textbf{第一轮 Prompt：}
		\begin{exampleblock}{}
			Role: 图像处理专家

			Task: Python+OpenCV实现去噪

			要求：高斯滤波、中值滤波对比
		\end{exampleblock}

		\column{0.5\textwidth}
		\textbf{第二轮 Prompt：}
		\begin{exampleblock}{}
			添加Docstring、解释适用场景、添加可视化、错误处理
		\end{exampleblock}
	\end{columns}

	\tiny
	\textbf{技巧：}多轮迭代、明确告诉AI需要改进什么
\end{frame}

% -----------------------------------------------------------------------------
% 3. 实战 2：透视变换
% -----------------------------------------------------------------------------

\begin{frame}{实战 2：透视变换}
	\tiny
	\textbf{挑战：}透视变换数学原理复杂，如何通过AI快速实现？

	\begin{columns}
		\column{0.5\textwidth}
		\textbf{第一轮：让AI解释原理}
		\begin{exampleblock}{}
			用简单比喻解释什么是透视变换
		\end{exampleblock}

		\textbf{AI比喻：}
		\begin{itemize}
			\item 透视变换 = 把斜片"拉正"
			\item getPerspectiveTransform = 计算"怎么拉"
			\item warpPerspective = 执行"拉"
		\end{itemize}

		\column{0.5\textwidth}
		\textbf{第二轮：实现代码}
		\begin{exampleblock}{}
			请根据上面的解释实现函数
		\end{exampleblock}

		\textbf{关键步骤：}预处理 $\to$ 找轮廓 $\to$ 定角点 $\to$ 透视变换
	\end{columns}
\end{frame}

% -----------------------------------------------------------------------------
% 4. 实战 3：自动识别答题卡边缘
% -----------------------------------------------------------------------------

\begin{frame}{实战 3：自动识别答题卡边缘（迭代优化）}
	\textbf{目标：} 通过 3 轮迭代 Prompt，让 AI 从写出简单 Canny 算子到写出鲁棒的轮廓提取逻辑。

	\vspace{0.3cm}

	\begin{columns}
		\column{0.5\textwidth}
		\textbf{第一轮：基础实现}

		\begin{alertblock}{基础 Prompt}
			请用 OpenCV 实现答题卡边缘检测。
		\end{alertblock}

		\textbf{AI 输出：}
		\begin{itemize}
			\item 简单的 Canny 边缘检测
			\item 查找轮廓
			\item 可能问题：噪点干扰、参数固定
		\end{itemize}

		\textbf{效果：}\textcolor{red}{\textbf{较差}}

		\column{0.5\textwidth}
		\textbf{第二轮：优化改进}

		\begin{exampleblock}{优化 Prompt}
			上面的代码检测效果不好，请改进：
			\begin{enumerate}
				\item 先进行图像预处理（去噪）
				\item 使用自适应阈值
				\item 筛选出最大的四边形
			\end{enumerate}
		\end{exampleblock}

		\textbf{AI 输出：}
		\begin{itemize}
			\item 高斯模糊去噪
			\item 自适应 Canny 阈值
			\item 面积筛选 + 四边形近似
		\end{itemize}

		\textbf{效果：}\textcolor{orange}{\textbf{良好}}
	\end{columns}
\end{frame}

\begin{frame}{实战 3（续）：第三轮迭代}
	\begin{columns}
		\column{0.5\textwidth}
		\textbf{第三轮：鲁棒完善}

		\begin{exampleblock}{完善 Prompt}
			请进一步优化，使其更鲁棒：
			\begin{enumerate}
				\item 处理极端拍摄角度
				\item 处理光照不均和阴影
				\item 添加透视变换矫正
				\item 封装成可复用的类
			\end{enumerate}
		\end{exampleblock}

		\textbf{AI 输出：}
		\begin{itemize}
			\item 完整的类封装
			\item CLAHE 预处理
			\item 透视变换矫正
			\item 调试可视化
		\end{itemize}

		\textbf{效果：}\textcolor{green}{\textbf{优秀}}

		\column{0.5\textwidth}
		\begin{block}{迭代优化的关键}
			\begin{itemize}
				\item \textbf{循序渐进}：每轮只解决1-2个问题
				\item \textbf{具体反馈}：告诉AI具体哪里不好
				\item \textbf{要求明确}：提出具体的改进要求
			\end{itemize}
		\end{block}
	\end{columns}
\end{frame}

% -----------------------------------------------------------------------------
% 5. 迭代优化总结
% -----------------------------------------------------------------------------

\begin{frame}{迭代优化演示：从简单到完善}
	\begin{center}
		\begin{tikzpicture}
			\node[draw, fill=red!20] (v1) at (0,0) {\shortstack{第1轮\\基础实现}};
			\node[draw, fill=yellow!20] (v2) at (4,0) {\shortstack{第2轮\\优化改进}};
			\node[draw, fill=green!20] (v3) at (8,0) {\shortstack{第3轮\\鲁棒完善}};

			\draw[->, thick] (v1) -- (v2) node[midway, above] {添加功能};
			\draw[->, thick] (v2) -- (v3) node[midway, above] {完善细节};
		\end{tikzpicture}
	\end{center}

	\vspace{0.3cm}

	\begin{columns}
		\column{0.33\textwidth}
		\textbf{第1轮：基础}
		\begin{itemize}
			\item 简单 Canny
			\item 查找轮廓
			\item 基本功能
		\end{itemize}
		\textbf{问题：}
		\begin{itemize}
			\item 噪点干扰
			\item 参数固定
			\item 不完整边缘
		\end{itemize}

		\column{0.33\textwidth}
		\textbf{第2轮：优化}
		\begin{itemize}
			\item 高斯模糊
			\item 自适应阈值
			\item 面积筛选
		\end{itemize}
		\textbf{改进：}
		\begin{itemize}
			\item 抗噪性增强
			\item 自动调参
			\item 鲁棒性提升
		\end{itemize}

		\column{0.33\textwidth}
		\textbf{第3轮：完善}
		\begin{itemize}
			\item 完整类封装
			\item CLAHE 预处理
			\item 透视变换
		\end{itemize}
		\textbf{特性：}
		\begin{itemize}
			\item 可复用性强
			\item 调试可视化
			\item 生产级可用
		\end{itemize}
	\end{columns}
\end{frame}

\begin{frame}{迭代优化的核心思想}
	\begin{block}{为什么需要多轮迭代？}
		\begin{itemize}
			\item \textbf{第1轮}：验证可行性，让代码"跑起来"
			\item \textbf{第2轮}：针对问题优化，让代码"好用"
			\item \textbf{第3轮}：完善细节，让代码"专业"
		\end{itemize}
	\end{block}

	\begin{alertblock}{关键技巧}
		\begin{itemize}
			\item 每轮只解决1-2个核心问题
			\item 不要一次性完成所有优化
			\item 每轮都要测试验证
		\end{itemize}
	\end{alertblock}

	\vspace{0.3cm}

	\begin{block}{与AI协作的迭代模式}
		\begin{enumerate}
			\item 第1轮Prompt：实现基础功能
			\item 第2轮Prompt：针对问题进行优化
			\item 第3轮Prompt：完善代码风格和注释
		\end{enumerate}
	\end{block}
\end{frame}

% -----------------------------------------------------------------------------
% 6. 实战任务：人脸检测
% -----------------------------------------------------------------------------

\begin{frame}[fragile]{实战任务：用AI辅助实现人脸检测}
	\textbf{Prompt示例：}

	\begin{exampleblock}{人脸检测 Prompt}
		请用Python和OpenCV实现一个人脸检测程序：

		功能：从图片中检测所有人脸，并用矩形框标注

		输入：图片文件路径

		输出：标注了人脸框的图片
	\end{exampleblock}

	\vspace{0.3cm}

	\begin{block}{具体要求}
		\begin{itemize}
			\item 使用OpenCV的Haar级联分类器
			\item 在每个人脸周围绘制绿色矩形框
			\item 显示检测到的人脸数量
			\item 代码有详细中文注释
		\end{itemize}
	\end{block}
\end{frame}

\begin{frame}{预期输出}
	\textbf{完整的可运行代码，包含：}
	\begin{enumerate}
		\item 详细的函数文档（Docstring）
		\item 中文注释
		\item 使用示例
		\item 参数说明
	\end{enumerate}

	\vspace{0.3cm}

	\begin{alertblock}{提示}
		如果AI生成的代码有问题，可以：
		\begin{itemize}
			\item 把错误信息发给AI，让它修复
			\item 询问AI如何优化代码
			\item 让AI解释代码的每一行
		\end{itemize}
	\end{alertblock}
\end{frame}

\begin{frame}{人脸检测代码示例}
	\begin{block}{核心代码}
		\begin{lstlisting}[language=Python, basicstyle=\ttfamily\tiny]
import cv2

def detect_faces(image_path):
    face_cascade = cv2.CascadeClassifier(
        cv2.data.haarcascades + 'haarcascade_frontalface_default.xml')
    img = cv2.imread(image_path)
    gray = cv2.cvtColor(img, cv2.COLOR_BGR2GRAY)
    faces = face_cascade.detectMultiScale(gray, 1.1, 5, (30,30))
    for (x, y, w, h) in faces:
        cv2.rectangle(img, (x,y), (x+w,y+h), (0,255,0), 2)
    return img
		\end{lstlisting}
	\end{block}

	\vspace{0.2cm}

	参数：\texttt{scaleFactor=1.1, minNeighbors=5, minSize=(30,30)}
	\\
	\tiny（注：参数需根据实际场景调整）
\end{frame}
